% Generated by Listas/make_statements.py from pratica-lecture-05.md, item 20. Do not edit.
\textbf{Source.} Mas-Colell, Whinston \& Green (1995), \emph{Microeconomic Theory}, Appendix A to Ch. 12 ``Infinitely Repeated Games and the Folk Theorem'', Exercise 12.AA.1\^{}A (in text), p.~420 {[}PDF 221{]}. Answer in the source: in the MWG solutions manual (which marks the printed statement as incorrect: first-printing erratum). Reproduced verbatim. (Context, not in the exercise: the stage game is a duopoly in which firm \(i\) chooses \(q_i\) and earns \(\pi_i(q_1, q_2)\); \(q^* = (q_1^*, q_2^*)\) is its Nash equilibrium; Nash reversion, Definition 12.AA.1 on p.~418, plays a prescribed outcome path until someone deviates and \(q^*\) forever after.)



\begin{quote}

\textbf{Exercise 12.AA.1:} Argue that no pair of actions \(q\) such that \(\pi_i(q_1, q_2) < \pi_i(q_1^*, q_2^*)\) for some \(i\) can be sustained as a stationary SPNE outcome path using Nash reversion.

\end{quote}



\textbf{Note (not part of the exercise).} The MWG solutions manual (p.~12-46) says the printed statement is incorrect (first-printing erratum), gives a counterexample, and states the closest correct claim: no \(q\) with \(\pi_i(q) < \pi_i(q^*)\) for \emph{all} \(i\) can be sustained by Nash reversion. The solution works both.

